\section{Pengukuran Energi dengan Ansatz \texttt{EfficientSU2}}
\texttt{EfficientSU2} didefinisikan sebagai salah satu jenis \textit{Hardware-Efficient Ansatz} (HEA) yang dirancang secara khusus untuk mengoptimalkan penggunaan gerbang kuantum pada perangkat keras era NISQ yang memiliki keterbatasan sumber daya. Arsitektur sirkuit ini terdiri dari lapisan gerbang rotasi tunggal dari grup $SU(2)$, seperti gerbang $R_y$ dan $R_z$, yang kemudian diikuti oleh lapisan gerbang keterbelitan (\textit{entanglement}) CNOT untuk menciptakan korelasi non-lokal antar qubit. Keunggulan utama dari desain \texttt{EfficientSU2} terletak pada fleksibilitasnya dalam mengatur jumlah pengulangan lapisan atau \textit{reps} guna menyeimbangkan antara tingkat ekspresivitas sirkuit dengan kerentanan terhadap akumulasi derau hardware. Dengan struktur yang sangat efisien ini, sirkuit kuantum mampu memetakan interaksi kompleks antar aset dalam Hamiltonian Ising tanpa memerlukan waktu koherensi yang lama atau jumlah gerbang yang melampaui kapasitas hardware saat ini.

Proses pengukuran energi dalam algoritma VQE dilakukan dengan mengevaluasi nilai ekspektasi dari operator Hamiltonian terhadap keadaan kuantum yang dihasilkan oleh sirkuit parameterisasi \texttt{EfficientSU2}. Mengingat Hamiltonian Ising sepenuhnya terdiri dari suku-suku operator \textit{Pauli-Z}, proses pengukuran dapat dilakukan secara langsung pada basis komputasi standar untuk mengekstraksi nilai eigen dari sistem fisik tersebut. Estimasi nilai ekspektasi energi diperoleh melalui perhitungan rata-rata statistik dari ribuan pengukuran (\textit{shots}) guna mereduksi kesalahan probabilitas yang melekat pada sifat stokastik dari perangkat keras kuantum. Strategi pemetaan parameter yang diterapkan dalam arsitektur \texttt{EfficientSU2} memastikan bahwa setiap perubahan pada variabel $\theta$ memiliki pengaruh yang terukur secara langsung terhadap perubahan total energi sistem. Dengan menggabungkan desain sirkuit yang efisien serta mekanisme pengukuran yang memiliki presisi tinggi, sistem kuantum mampu mengekstrak solusi portofolio optimal dari ribuan kemungkinan konfigurasi dalam ruang pencarian.
